\documentclass[11pt]{article}
\newcommand{\LabNumber}{2}
\newcommand{\LabTitle}{Arrays and ArrayLists}
\newcommand{\LabTopic}{Fixed arrays, insertion/deletion, dynamic resizing, amortised analysis}
\input{common.tex}

\begin{document}
\labheader

\section*{Objectives}
\begin{itemize}[leftmargin=*,nosep]
  \item Declare, initialise, and traverse one-dimensional arrays in Java.
  \item Implement insertion and deletion in a sorted array and understand the $O(n)$ shift cost.
  \item Build a \texttt{Scoreboard} class that maintains a sorted top-$k$ array.
  \item Implement a generic \texttt{ArrayList<E>} backed by an \texttt{Object[]} that doubles when full.
  \item Analyse the amortised $O(1)$ cost of \texttt{add} using the accounting method.
\end{itemize}

\begin{summary}[Quick Reference -- Arrays]
\textbf{Declaring arrays.}
\begin{lstlisting}
int[] a = new int[5];                      // all zeros
String[] names = {"Alice", "Bob", "Carol"};
int len = a.length;                        // field, NOT a method
\end{lstlisting}
\textbf{Shifting right} to insert at index $k$ (loop \emph{backward}):
\begin{lstlisting}
for (int i = last; i >= k; i--) a[i+1] = a[i];
a[k] = newValue;
\end{lstlisting}
\textbf{Complexity.}\quad Random access: $O(1)$. Insert/delete (keeping order): $O(n)$.
\end{summary}

\begin{summary}[Quick Reference -- ArrayList]
\textbf{Resizing strategy.}\quad When full, allocate double the capacity and copy all elements.
\begin{lstlisting}
private void resize() {
    Object[] bigger = new Object[data.length * 2];
    for (int i = 0; i < size; i++) bigger[i] = data[i];
    data = bigger;
}
\end{lstlisting}
\textbf{Amortised analysis.}\quad Starting from capacity 1, resizes at sizes $1,2,4,\ldots,n$
cost $1+2+4+\cdots+n = 2n-1$ total copies for $n$ adds: $O(1)$ amortised per \texttt{add}.\\
\textbf{Key operations.}
\begin{tabular}{@{}lll@{}}
\toprule
\texttt{get(i)}/\texttt{set(i,e)} & $O(1)$ & always \\
\texttt{add(e)} at tail & $O(1)$ & amortised \\
\texttt{add(i,e)} at index & $O(n)$ & shift cost \\
\texttt{remove(i)} & $O(n)$ & shift cost \\
\bottomrule
\end{tabular}
\end{summary}

\subsection*{Quick Experiments (Try It)}

\begin{tryit}{Off-by-one index}
\begin{lstlisting}
int[] a = {10, 20, 30, 40, 50};
for (int i = 0; i <= a.length; i++)   // deliberate bug
    System.out.println(a[i]);
\end{lstlisting}
What exception is thrown and on which iteration? Fix the loop bound.
\end{tryit}

\begin{tryit}{Watch capacity double}
Add a \texttt{capacity()} accessor to your \texttt{ArrayList}.
Start with capacity 2 and append elements one by one; print the capacity after each add.
Count the total number of element copies across 16 appends. Is it less than $2 \times 16$?
\end{tryit}

\section*{In-Lab Exercises (target: 2 hours)}

\begin{labexercise}{Array Statistics \hfill {\small \itshape (20 min)}}
Read $n$ integers from the user (first read $n$, then the values).
Print: the minimum, maximum, and average (as \texttt{double}).
Handle $n = 1$ correctly.
\end{labexercise}

\begin{labexercise}{\texttt{Scoreboard} Class \hfill {\small \itshape (40 min)}}
Implement \texttt{GameEntry} (name + score, \texttt{toString()}) and
\texttt{Scoreboard(int capacity)} that keeps at most \texttt{capacity} entries in descending score order.
\begin{itemize}[leftmargin=*,nosep]
  \item \texttt{add(GameEntry e)}: insert if the board is not full or \texttt{e} beats the lowest score; shift as needed.
  \item \texttt{remove(int i)}: remove entry at index \texttt{i} and shift left.
  \item \texttt{toString()}: print all entries, one per line.
\end{itemize}
Test with at least 6 entries and one removal.
\end{labexercise}

\begin{labexercise}{Implement \texttt{ArrayList<E>} \hfill {\small \itshape (50 min)}}
Build the class with an initial capacity of 4:
\begin{itemize}[leftmargin=*,nosep]
  \item \texttt{int size()}, \texttt{boolean isEmpty()}, \texttt{int capacity()}.
  \item \texttt{E get(int i)} --- throws \texttt{IndexOutOfBoundsException} if out of range.
  \item \texttt{void set(int i, E e)}.
  \item \texttt{void add(E e)} --- appends; doubles capacity when full.
  \item \texttt{void add(int i, E e)} --- inserts at index \texttt{i}, shifting right.
  \item \texttt{E remove(int i)} --- removes and returns; shifts left.
  \item Shrink policy: if after removal $\texttt{size} < \texttt{capacity}/4$, halve capacity (min 4).
  \item \texttt{String toString()} --- e.g.\ \texttt{[A, B, C] (size=3, cap=4)}.
\end{itemize}
Test: append 6 elements, insert at index 2, remove index 0, print. Verify one resize occurred.
\end{labexercise}

\begin{aicollab}
\textbf{Suggested prompts:}
\begin{itemize}[leftmargin=*,nosep]
  \item ``Why must I loop \emph{backwards} when shifting array elements right for insertion?''
  \item ``Explain amortised $O(1)$ for \texttt{ArrayList.add} using the accounting/potential method.''
\end{itemize}
\medskip\textbf{Verify:} Ask the AI to predict the total element copies for 32 appends starting from
capacity 1 with a doubling strategy. Compute it yourself and compare.
\end{aicollab}

\takehomebanner

\begin{trace}[Insertion shift]
Trace the following code for \texttt{a = \{2, 5, 8, 0\}}, \texttt{size = 3}, \texttt{v = 6}.
Show the array after each loop iteration and the final state.
\begin{lstlisting}
int pos = size;
while (pos > 0 && a[pos-1] > v) { a[pos] = a[pos-1]; pos--; }
a[pos] = v;
\end{lstlisting}
\end{trace}

\begin{trace}[ArrayList capacity trace]
Starting from capacity 2, trace \texttt{capacity()} after each call (use shrink-at-$<1/4$ policy, min 4):
\texttt{add(A), add(B), add(C), add(D), add(E), remove(0), remove(0), remove(0)}.
\end{trace}

\begin{writecode}[Binary search on a sorted array]
Write \texttt{int binarySearch(int[] a, int n, int target)}: returns the index of \texttt{target}
in the sorted array of length \texttt{n}, or $-1$ if absent. Use the iterative approach.
\end{writecode}

\begin{drawex}{Amortised accounting table}
Draw a table: add \#, element, capacity before, resize? (Y/N), copies made.
Fill it for adds 1--16 starting from capacity 1.
Sum the copies and show the total is less than $2 \times 16 = 32$.
\end{drawex}

\vspace{4pt}
\noindent\textbf{Submission.} Save files as \texttt{lab2\_ex*.java} and submit on the course site.

\end{document}
